\section{Schematic of the NREM/REM Cycle}\label{sec:nrem_rem_schematic}

\textit{[Placeholder: Insert short refs for stage neuromodulation profiles (e.g., noradrenergic/cholinergic 'amnesic' state of REM), rodent/human replay studies, and evidence for REM-linked generative/associative search.]}

\begin{figure}[htbp]
  \centering
  \framebox{\parbox{0.8\textwidth}{\centering
    \vspace{5cm}
    \textbf{Placeholder: NREM/REM Learning Schematic} \\
    \small\textit{A diagram showing NREM (replay/consolidation) and REM (generalization/augmentation) as two parts of a learning cycle.}
    \vspace{5cm}
  }}
  \caption{Schematic of the proposed two-stage learning algorithm. (1) NREM/SWS facilitates replay and consolidation of specific episodic memories. (2) REM sleep injects "adversarial" or "augmented" data (bizarre dreams), forcing the system to generalize its new knowledge, thus preventing overfitting.}
  \label{fig:schematic}
\end{figure}
